\documentclass[11pt]{article}
\usepackage[margin=1in]{geometry}
\usepackage{amsmath}
\usepackage{amssymb}
\usepackage{booktabs}
\usepackage{listings}
\usepackage{xcolor}
\usepackage{hyperref}

\lstset{basicstyle=\ttfamily\small,breaklines=true,frame=single,
        backgroundcolor=\color{gray!8},columns=fullflexible}
\newcommand{\code}[1]{\texttt{#1}}

\title{The reversed-wind advection adjoint in GCHP:\\
       \code{reverse\_wind\_pure} versus \code{reverse\_wind}}
\author{}
\date{}

\begin{document}
\maketitle

\noindent\emph{This note documents the difference between the two reversed-wind
advection-adjoint variants selected at run time by the \code{ADV\_ADJ\_METHOD}
key in \code{GCHP.rc}. Section~\ref{sec:detailed} is a detailed technical
description written to be self-contained for a reader who is not an expert in
finite-volume advection. Section~\ref{sec:manuscript} is a condensed,
manuscript-ready paragraph. All source references are to the
\code{GCHP\_fix\_adv} tree (module
\code{FVdycoreCubed\_GridComp/fvdycore/model/tp\_core.F90} and
\code{fv\_tracer2d.F90}).}

%======================================================================
\section{Detailed technical description}
\label{sec:detailed}

\subsection{The run-time switch}

The advection adjoint is selected by a single key in \code{GCHP.rc},
\begin{lstlisting}
ADV_ADJ_METHOD: reverse_wind
\end{lstlisting}
which admits three values: \code{reverse\_wind\_pure}, \code{reverse\_wind}, and
\code{exact\_linear}. The first two are the subject of this note; the third is a
true discrete transpose available only for the limiter-free \code{hord=2}/
\code{kord=17} path and is not discussed further here.

\subsection{A minimal primer on the forward transport scheme}

Tracer mass (here CO$_2$) is transported by the wind with GEOS-Chem/FV3's
finite-volume, flux-form advection on the cubed-sphere grid. On each dynamics
step the scheme performs, for every grid cell, three operations:
\begin{enumerate}
  \item \textbf{Sub-grid reconstruction.} The tracer field is represented inside
  each cell not by a single flat value but by a piecewise-parabolic profile
  (the Piecewise Parabolic Method, PPM), built from a cell and its neighbours in
  the routines \code{xppm} (east--west) and \code{yppm} (north--south).
  \item \textbf{Face fluxes.} The mass crossing each cell face is obtained by
  integrating that reconstructed profile over the region swept by the flow in
  one time step. The flow enters through the \emph{Courant number}
  $c=u\,\Delta t/\Delta x$, and the scheme is \emph{upwind-biased}: the tracer
  that crosses a face is taken from the cell on the side the wind is blowing
  \emph{from}.
  \item \textbf{Cell update.} Each cell's mass is updated as (flux in)
  $-$ (flux out).
\end{enumerate}

Two additional features matter for the adjoint, and both live at the
\emph{seams} of the cube --- the edges and corners where the six faces meet and
the coordinate lines are discontinuous:
\begin{itemize}
  \item \code{copy\_corners} copies tracer values from a neighbouring face into
  the ``ghost'' halo cells surrounding each of the eight cube corners, so that
  the reconstruction stencil has data to work with there;
  \item at each face \emph{edge}, \code{xppm}/\code{yppm} use special one-sided
  boundary stencils that blend an off-face ghost value with interior values to
  build the edge parabola.
\end{itemize}

\subsection{The reversed-wind approximation}

The adjoint model must apply the mathematical \emph{transpose} of the entire
forward transport operator, propagating sensitivities ($\partial J/\partial$
state) backward in time. Forming that transpose exactly is expensive and, for
the flux/slope limiters used in production (\code{hord\_tr=12}, \code{kord\_tr=8}),
has not been implemented; that exact path exists only for the limiter-free case
(\code{exact\_linear}).

Both \code{reverse\_wind\_pure} and \code{reverse\_wind} instead use a cheap
approximation: \textbf{they re-run the ordinary forward transport code with the
mass fluxes (winds) negated}. Reversing the wind approximates the transpose
because advection is close to a simple upwind operator; the approximation is
\emph{exact} when the flow is spatially uniform and degrades as the flow
develops shear and curvature. Crucially, both variants execute the identical
call chain
\code{tracer\_2d}~$\rightarrow$~\code{fv\_tp\_2d}~$\rightarrow$~
\code{xppm}/\code{yppm} with negated fluxes; \emph{neither} forms the exact
transpose.

\subsection{The single distinguishing flag}

The only thing that differs between the two variants is one module-level logical,
set in \code{fv\_tracer2d.F90}:
\begin{lstlisting}
tp_core_is_adjoint = l_isAdjoint .and. (trim(adv_adj_method) == 'reverse_wind')
\end{lstlisting}
Thus
\begin{itemize}
  \item \code{reverse\_wind\_pure} $\Rightarrow$
        \code{tp\_core\_is\_adjoint = .false.}: the winds are negated and the
        forward operator is run verbatim, with \emph{no} further modification;
  \item \code{reverse\_wind} $\Rightarrow$
        \code{tp\_core\_is\_adjoint = .true.}: three direction-aware corrections
        are switched on at the cube seams.
\end{itemize}

\subsection{Why corrections are needed, and what they are}

When the wind is negated, the \emph{upwind direction flips}, but the forward
code's corner-copy and edge stencils were written for the \emph{forward} upwind
direction. Consequently, at every face edge and corner they reach for the wrong
neighbouring cells. The three corrections repair exactly those locations.

\paragraph{Correction 1 --- the corner copy becomes its transpose.}
(\code{tp\_core.F90}, in \code{fv\_tp\_2d}.) Where the pure variant calls
\code{copy\_corners} (which \emph{scatters} a corner tracer value \emph{out} to
the ghost halo), the corrected variant calls \code{copy\_corners\_ad}, the
adjoint of that copy. Instead of scattering, it \emph{accumulates} the ghost-cell
sensitivities \emph{back into} the source real cells and then zeroes the ghosts:
\begin{lstlisting}
! Adjoint of copy_corners: accumulate ghost corner sensitivities back to
! the source real cells and zero the ghost corners.
q(j,1-i) = q(j,1-i) + q(i,j)
q(i,j)   = 0.
\end{lstlisting}
This is the correct way to carry sensitivity across a cube corner when the
integration runs backward.

\paragraph{Correction 2 --- edge reconstruction stencils switch to interior-only
upwind.} (\code{xppm} and \code{yppm}.) At a face edge the forward PPM builds the
edge parabola from a \emph{centred} blend of an off-face ghost value and interior
values. The correction detects, from the sign of the (reversed) Courant number,
when the adjoint flow points \emph{outward} across that edge --- $c<0$ at the low
edge, $c>0$ at the high edge --- and in that case replaces the centred blend with
an \emph{interior-only} reconstruction:
\begin{lstlisting}
! Outward adjoint flow (c<0): upwind is interior -- use interior-only reconstruction.
if (tp_core_is_adjoint .and. c(is,j) < 0.) &
   al(1) = ((2.*dxa(1,j)+dxa(2,j))*q1(1)-dxa(1,j)*q1(2))/(dxa(1,j)+dxa(2,j))
\end{lstlisting}
Because the flow is reversed, the physically upwind cells are now the interior
ones rather than the ghosts, so the stencil stops reaching across the seam. The
same test is applied at all four face edges (\code{is==1} and
\code{(ie+1)==npx} in $x$; \code{js} and \code{(je+1)==npy} in $y$).

\paragraph{Correction 3 --- the same treatment for the parabola-edge slopes.}
(The \code{bl}/\code{br} boundary blocks of \code{yppm}, and the corresponding
$x$ blocks.) The one-sided parabola end-slopes at the face edges are likewise
rebuilt with an interior-only upwind stencil (with the monotonicity bounds
re-applied to the interior triplet) wherever the reversed flow is outward. This
is the slope-limited counterpart of Correction~2.

\subsection{What the corrections do \emph{not} fix}

The three corrections act only at the cube \emph{seams}. They do not alter the
\emph{interior} PPM operator, for which the reversed-wind sweep is still not the
exact transpose: the cross-derivative (inner, cross-swept) step of
\code{fv\_tp\_2d} divides the intermediate half-updated tracer by the swept-area
factors \code{ra\_x}/\code{ra\_y} in a way that a faithful transpose would not.
This residual, wind-gradient--dependent term is common to both reversed-wind
variants and is removed only by the exact transpose (\code{exact\_linear}).

\subsection{Empirical consequence}

We validated both variants against finite-difference (FD) gradients on the
native C24 cubed-sphere grid in a one-month, advection-only experiment
(\code{exp3\_adv\_only}: transport-only physics, control $=$ a per-cell
multiplicative flux scale $\sigma$, cost $J$ a quadratic XCO$_2$ misfit). The FD
reference uses central differences $\big(J(\sigma{+}\varepsilon)-J(\sigma{-}
\varepsilon)\big)/2\varepsilon$ at $\varepsilon=0.05$; the forward model is
bit-identical across the adjoint-only variants, so a single set of FD runs
serves both. Results at the sampled cells, including three large-signal
``hotspots'' (Ethiopia, Colombia, Brazil):

\begin{center}
\begin{tabular}{lccc}
\toprule
variant & correlation with FD & through-origin slope & hotspot $g_{\rm adj}/g_{\rm FD}$ \\
\midrule
\code{reverse\_wind\_pure} & $-0.886$ & $-6.49$ & $-4.4,\ -8.2,\ -2.9$ \\
\code{reverse\_wind}       & $+0.988$ & $+4.10$ & $+4.0,\ +4.4,\ +3.5$ \\
\bottomrule
\end{tabular}
\end{center}

The uncorrected variant is \emph{anti-correlated} with the truth: with the cube
seams untreated, the corner copies and outward-edge stencils point the wrong way
and the recovered gradient has the wrong sign at essentially every cell. Enabling
the three corrections restores the correct sign and a strong correlation
($+0.99$), leaving a near-constant amplitude bias of roughly a factor of four
that is attributable to the un-transposed interior cross-derivative term and is
removed only by the exact transpose.

%======================================================================
\section{Manuscript paragraph}
\label{sec:manuscript}

\noindent\emph{Condensed version for the paper. Numbers correspond to the
\code{exp3\_adv\_only}, one-month, native-C24 validation; adjust the
configuration clause to match the manuscript's notation.}

\bigskip

The advection adjoint is evaluated with a reversed-wind approximation: rather
than forming the exact discrete transpose of the piecewise-parabolic (PPM)
cubed-sphere transport operator---which, for the flux-limited scheme used here,
is not tractable---the ordinary forward advection is re-integrated with the mass
fluxes negated. This approximation is exact for spatially uniform flow and
degrades with flow shear and curvature. We compare two variants. In the
uncorrected variant the winds are simply negated and the forward operator is
applied verbatim; because reversing the flow inverts the upwind direction while
the operator's cube-seam bookkeeping (the corner halo copy and the one-sided
face-edge reconstruction stencils) remains coded for the forward direction, the
recovered sensitivities are \emph{anti-correlated} with a finite-difference
reference (correlation $-0.89$) and carry the wrong sign at essentially every
grid cell. In the corrected variant three direction-aware modifications are
applied at the cube seams: (i) the inter-face corner copy is replaced by its
adjoint, which accumulates ghost-cell sensitivities back onto the source cells
rather than scattering them outward; and (ii,iii) at each face edge, wherever the
reversed flow is outward (diagnosed from the sign of the Courant number), the
centred ghost-plus-interior PPM edge reconstruction and its limited end-slopes
are replaced by an interior-only upwind reconstruction, so the stencil no longer
reaches across the seam in the wrong direction. With these corrections the
adjoint sensitivities are strongly and correctly correlated with finite
differences (correlation $+0.99$, correct sign at all validated cells). A
residual, near-constant amplitude bias of approximately a factor of four remains;
it originates in the interior cross-derivative step of the two-dimensional PPM
operator, which the reversed-wind construction does not transpose exactly, and is
eliminated only by the exact discrete adjoint (available for the limiter-free
scheme). For the gradient-based optimisation used here the corrected reversed-
wind adjoint provides correctly signed, spatially faithful sensitivities at
low cost, with the residual amplitude bias absorbed into the step-length control
of the minimiser.

\medskip

\noindent\textbf{One-sentence summary.}
Negating the winds alone yields sign-reversed advection sensitivities
($r=-0.89$); additionally replacing the cube-corner halo copy with its adjoint
and switching the face-edge PPM stencils to interior-only upwind wherever the
reversed flow is outward makes the reversed-wind adjoint correctly signed and
strongly correlated with finite differences ($r=+0.99$), leaving only a
$\sim$4$\times$ amplitude bias from the un-transposed interior cross-derivative
term.

\end{document}
